% ============================================================
% Capítulo 2: Planteamiento del problema
% ============================================================
\chapter{Planteamiento del Problema}
\label{cap:problema}

\section{Descripción del problema}

Exposición clara de la situación problemática que motiva la investigación.

\section{Formulación del problema}

\subsection{Pregunta general}
¿Cuál es la pregunta central que guía la investigación?

\subsection{Preguntas específicas}
\begin{itemize}
  \item Pregunta específica 1.
  \item Pregunta específica 2.
  \item Pregunta específica 3.
\end{itemize}

\section{Objetivos}

\subsection{Objetivo general}
Enunciar el objetivo general (verbo en infinitivo).

\subsection{Objetivos específicos}
\begin{itemize}
  \item Objetivo específico 1.
  \item Objetivo específico 2.
  \item Objetivo específico 3.
\end{itemize}

\section{Justificación}

Razones que fundamentan la importancia y pertinencia del estudio (teórica,
práctica, social, tecnológica).

\section{Alcance y limitaciones}

Delimitación del estudio (temporal, espacial, temática) y restricciones
encontradas.
